\scanpage{25}
% lines 1–6
à une forme quadratique non dég -- qu'on peut dès lors, pour une base convenable, mettre sous la forme $xy+z^2$. Quitte à multiplier $\omega$ par un scalaire (ce qui est loisible) on peut supposer que $\omega=e_1\wedge e_2\wedge e_3$, ce qui nous ramène au cas déjà traité.

% lines 7–9
Lorsque $\varphi$ est associé à une forme quadratique $q$, on tire de la deuxième formule (74), compte tenu de $\varphi(x,x)=2q(x)$ :

% lines 10
\[
\operatorname{Cas}_E(x,x)=-2\Delta\varphi(x,x)=-4\Delta q(x) \tag*{(75)}
\]

% lines 11–12
ce qui, dans le cas d'une forme \inserted{$q$} spéciale orthogonale pour $\omega$, d'où $\Delta=-2$, donne

% lines 13
\[
\operatorname{Cas}_E(x,x)=8q(x) \tag*{(76)}
\]

% lines 14–17
Donc dans ce cas, \deleted{\unsure{on a}} si 2 est inv, (ou plus gén. non div. de 0) $q(x)$ se \uline{déduit} de la structure d'algèbre de Lie sur $E$, comme l'\unsure{op.|exp.|opér.} de Casimir divisé par $8\ldots$

% lines 18–19
On \unsure{tire|tirera|trouve} de =, par la première formule

% lines 20
\[
\operatorname{Cas}_{E'}(x',x')=2\varphi'(x',x')=2q'(x') \tag*{\deleted{\unsure{(77)}}(77)}
\]

% lines 21–24
i.e. on retrouve la forme quadratique $q'(x')$ (ou la forme bil $\varphi'$ dont elle provient) à l'aide de la structure d'Algèbre de Lie de $E'_{\deleted{\unsure{\mathfrak g}}}=\mathfrak g$.

% lines foot
% archive sheet number: 24
